\documentclass[a4paper,11pt]{ltjsarticle}
%\documentclass[a4paper,12pt]{jsarticle}

%\usepackage[kozuka-pr6n]{pxchfon}

\usepackage[no-math,deluxe,expert,
hiragino-pron,
%hiragino-pro,
%morisawa-pro,
%yu-win,
%ms-hg
%kozuka-pr6n
]{luatexja-preset}
%\usepackage[utf8]{inputenc} 

\usepackage{arydshln,mleftright,}


\usepackage{academicons}
%\usepackage{mathfont}
%\setmainfont
%{PalladioPro}
%{adobegaramondpro}
%{CalistoMTStd}
%\setmainfont[Scale=0.97]
%%{CaeciliaLTStd}
%{PalladioPro}
%{AdobeTextPro}
%{Livingston}
%Livingston-Italic.otf
%Livingston-Medium.otf
%Livingston-MediumItalic.otf

%\setmainfont{adobetextpro}

%\mathfont[symbols]{PalladioPro}
%\mathfont[greekupper,greeklower]{PalladioPro}
%\mathfont{PalladioPro}



%\setsansfont[Ligatures=TeX]{TeXGyreHeros}
%\setmainjfont
%%[
%%BoldFont=
%%%%A-OTF-MidashiMinPro-MA31
%%%%A-OTF-RyuminPro-ExBold
%%%%HiraginoMinchoProN-W6
%%A-OTF-ReimYonzPro-Bold
%%]
%%[BoldFont=A-OTF-MidashiMinPro-MA31]w
%%[BoldFont=A-OTF-RyuminPro-Bold]
%%[BoldFont=A-OTF-UDReiminPro-Bold]
%%[BoldFont=TPMinStdN-CLoM]
%[BoldFont=HiraginoMinchoStdNW6]
%{
%%A-OTF-FutoMinA101Pr5-Bold
%%A-OTF-RyuminPr5-Regular
%%HamaMinStdN-TextR
%%A-OTF-ReimPr6-Regular
%HiraginoMinchoStdNW3
%%A-OTF-UDReiminPro-Light
%%A-OTF-ReimYonzPro-Regular
%%TPMinStdN-CLoR
%%HiraginoUD Std W4
%%HackGen35Nerd-Regular
%%HiraginoUDMinchoW4
%%UDDigiKyokashoStdN-Regular
%%Hiragino Mincho StdN
%%A-OTF-MidMiMA1Std-Bold
%%KozMinPr6N-Medium
%%A-OTF-FutoMinA101Pro-Bold
%%A-OTF-RyuminPro-Regular
%}
\setsansjfont
%[
%BoldFont=AxisBasicStdN-M
%%%A-OTF-MidashiMinPro-MA31
%%%A-OTF-RyuminPro-ExBold
%%%HiraginoMinchoProN-W6
%%A-OTF-ShinGoPro-Bold.otf
%]
%[BoldFont=KozGoPr6N-Bold]{A-OTF-GothicMB101Pro-Reg}
[BoldFont=A-OTF-GothicMB101Pro-DeBold]{A-OTF-GothicMB101Pro-Reg}
%[BoldFont=hiraginoUDMaruGoStdW6]{hiraginoUDMaruGoStdW4}
%[BoldFont=TPGothicStdN-M]{TPGothicStdN}
%{A-OTF-UDShinMGoPro-Regular}
%%\newjfontfamily\jisninety[CJKShape=JIS1990]{HaranoAjiMincho-Regular}%\usepackage{mymtpro2,myryumin}
%


%\setmathfont{asanamath}




\usepackage{bbold}
%\usepackage{mathpazo,newpxmath}
\usepackage[
palatino
%times
%palatino
]{fontsetup}
%\setmainfont{AdobeTextPro}
%\setsansfont[Ligatures=TeX]
%{FrutigerNeue1450Pro}
%{Humanist777bt}
\setsansfont{Humanist777bt}



\usepackage[margin=12mm]{geometry}

\usepackage{paralist}

\title{研究室運営方針}
\author{入谷亮介}
\begin{document}
\maketitle
\section{学術的なミッション・ビジョン}
当ラボが目指す研究の方向性と、共有すべき価値観です。

\begin{compactitem}
    \item \textbf{知的好奇心の追求:} 既存の枠組みにとらわれず、本質的な問いに向き合うことを最優先します。
    \item \textbf{質の重視:} 論文の数よりも、一つ一つの研究の質とインパクト、結果を洗練された形でプレゼンすること、そして論理の強固さを重視します。
    \item \textbf{異分野融合と協働:} 異なるバックグラウンドを持つ研究者との対話を歓迎し、新しい視点を取り入れることを恐れず、また、お互いのバックグラウンドの違いを尊重します。単純な上下関係や順序関係ではなく、お互いから学び合えることを見つけ、研究や知的活動につなげます。
    \item \textbf{社会への還元:} 基礎研究の深堀りを通じて、長期的には社会や人類の知識基盤に貢献することを目指します。
\end{compactitem}

\section{行動規範 (Code of Conduct)}
メンバー全員が安心して研究に打ち込めるよう、以下の規範を遵守します。

\begin{compactitem}
    \item \textbf{相互尊重:} 互いの専門性、背景、信条を尊重します。あらゆるハラスメント（アカデミック、セクシャル、パワーハラスメント等）を一切許容しません。本項の適用は研究室内部に限らず、研究室外活動でも同様です。
    \item \textbf{研究公正 (Research Integrity):} データの捏造、改ざん、盗用は行いません。誠実なデータ管理と報告を徹底します。
    \item \textbf{心理的安全性:} 議論における反論は「人格への攻撃」ではなく「真理への共同作業」です。わからないことを「わからない」と言える、失敗を隠さずに報告できる環境を作ります。
    \item \textbf{責任ある行動:} ラボのメンバーとしての自覚を持ち、共有スペースや機材を大切に扱います。
\end{compactitem}

\section{教育・研究環境 (Education \& Research Environment)}
PIと学生・スタッフの関係性、および成長のための指針です。

\begin{compactitem}
    \item \textbf{自律性の尊重:} 学生は「労働力」ではなく「独立した研究者の卵」として扱います。手取り足取り教えるのではなく、自ら考え、決定し、行動することを強く推奨します。
    \item \textbf{失敗への寛容さ:} 挑戦した結果の失敗は歓迎します。失敗から何を学び、次はどうするかを議論することに時間を割きます。
    \item \textbf{キャリア支援:} アカデミア、企業、官公庁など、多様なキャリアパスを等しく尊重し、支援します。
    \item \textbf{オープンな議論:} 定期的なミーティングに加え、Slackや対面での非公式な議論を推奨します。PIへの連絡は遠慮なく行ってください。
\end{compactitem}

\section{運営・実務ルール (Operations \& Logistics)}
日々のラボ運営に関する具体的な取り決めです。

\begin{compactitem}
    \item \textbf{勤務時間とコアタイム:} 原則として裁量労働制を重視しますが、コミュニケーション確保のため、平日の日中（例: 11:00--16:00）は連絡がつく状態を推奨します。
    \item \textbf{コミュニケーションツール:} 日常の連絡は Slack (または Teams) を使用します。緊急時を除き、深夜や休日の即時返信を求めることはありません。
    \item \textbf{データ管理:} 研究データはラボ指定のサーバー/クラウドに定期的にバックアップしてください。ハードディスクの紛失・破損リスクに常に備えてください。
    \item \textbf{ラボ当番:} 共有スペースの清掃、セミナーの準備などはメンバーで分担して行います。
    \item 半年に一度を目安に、1対1で30分ほど話す面談時間を設けましょう。もちろん、キャリア・研究・勉学・その他についての相談はいつでも受け付けます。
\end{compactitem}

\section{PIの取扱説明書 (User Guide to PI)}
円滑なコミュニケーションのための、PI（主宰者）に関するメモです。

\begin{compactitem}
%    \item メールの返信が24時間以内にない場合は、見落としている可能性が高いため、遠慮なく再送してください。
    \item 研究のアイデア出しや、数式の変形についての議論はいつでも歓迎します。ドアが空いている時は、アポイントメントなしで居室を訪ねても構いません。
    \item 私は愛想笑いをあまりしませんが、それは怒っているということではありません。真面目に聞いているときです。
    \item 私の質問は質問以外の意味を持ちません。質問による指摘（反語疑問文の利用）をしません。
    \item ミーティングを延期する必要などが生じた場合は、キャンセルした側が、具体的なりスケジュールを提案してください。
\end{compactitem}


\section{その他、細かいこと：健全な研究環境を維持するために}
	\begin{compactitem}
		\item デスクやオフィスは、構成員（含、PI）が大学に借りているものであり、プライベートではなく共有スペースです。ゴミは適切に片付け、衛生的に保ちます。冷蔵庫に物をいれる場合は、入れた日付と名前を必ず書いてください。未記入のものや、古くなったものは、こちらの判断で（なるべく予告のうえで）処分します。処分されて困るものは入れないでください。
		\item 研究室に到着して人に会ったらできるだけ挨拶をしましょう。PIに挨拶しにいく必要はありません。帰る時も同様です。
		\item 研究室に来なくなってしまった学生に気づいた場合は、なるべくPIに報告してください。
		\item 敷地内での飲み会やパーティなどの開催は、大学・大学院のルールに準拠します。
		\item 来客時、懇親会を開催することがあります。お店を探したりするのが得意な方にはお手伝いを依頼することがありますが、それは参加を強要するものではありませんので、美味しいお店を教えて下さい。
		\item あだ名でお互いに呼び合うことは構いませんが、行動規範項目に目を通し、適切な呼び方を心がけてください。私は男女問わず「さん」付けで呼びます。
	\end{compactitem}
\end{document}